\documentclass[11pt,a4paper]{article}

%% ── packages ─────────────────────────────────────────────────────────────────
\usepackage{mathptmx}
\usepackage[T1]{fontenc}
\usepackage[top=2.54cm, bottom=2.54cm, left=2.54cm, right=2.54cm]{geometry}
\usepackage{booktabs}
\usepackage{tabularx}
\usepackage{array}
\usepackage{float}          % [H] placement — table pinned exactly here
\usepackage{microtype}
\usepackage{fancyhdr}
\usepackage{titlesec}
\usepackage{ragged2e}
\usepackage{amsmath}
\usepackage{parskip}
\usepackage{hyperref}
\usepackage{caption}

%% ── Roman numeral table counter ──────────────────────────────────────────────
%% This makes \thetable output I, II, III ... so both captions and \ref are Roman
\renewcommand{\thetable}{\arabic{table}}

%% ── hyperref ─────────────────────────────────────────────────────────────────
\hypersetup{colorlinks=true, linkcolor=black, urlcolor=black, citecolor=black,
  pdftitle={Research Data Codebook - HIV and TB Kenya},
  pdfauthor={Otiende, Achia, Mwambi}}

%% ── IEEE-style table captions ────────────────────────────────────────────────
\captionsetup[table]{
  labelfont     = {sc,bf,small},
  textfont      = {sc,small},
  labelsep      = newline,
  skip          = 4pt,
  position      = above,
  justification = centering,
  singlelinecheck = false
}

%% ── column types ─────────────────────────────────────────────────────────────
\newcolumntype{L}[1]{>{\RaggedRight\arraybackslash}p{#1}}
\newcolumntype{C}[1]{>{\centering\arraybackslash}p{#1}}

%% ── inline code ──────────────────────────────────────────────────────────────
\newcommand{\code}[1]{\texttt{\small #1}}

%% ── section headings — IEEE roman-numeral centered small-caps ─────────────────
\setcounter{secnumdepth}{2}
\renewcommand\thesection{\arabic{section}}
\renewcommand\thesubsection{\arabic{subsection}}

\titleformat{\section}[block]
  {\normalsize\bfseries\scshape\centering}
  {\thesection.}{0.4em}{}
\titleformat{\subsection}[block]
  {\normalsize\itshape}
  {\thesubsection.}{0.4em}{}
\titlespacing*{\section}{0pt}{14pt}{6pt}
\titlespacing*{\subsection}{0pt}{10pt}{4pt}

%% ── header / footer ──────────────────────────────────────────────────────────
\pagestyle{fancy}\fancyhf{}
\renewcommand{\headrulewidth}{0.4pt}
\fancyhead[C]{\small\itshape
  Codebook --- HIV \& TB Joint Spatio-Temporal Analysis, Kenya (2012--2017)}
\fancyfoot[C]{\small\thepage}

%% ─────────────────────────────────────────────────────────────────────────────
\begin{document}
\setlength{\parindent}{1em}
\setlength{\parskip}{4pt}

%% ── TITLE BLOCK ──────────────────────────────────────────────────────────────
\begin{center}
{\Large\bfseries
  Research Data Codebook: Bayesian Joint Spatio-Temporal\\
  Modelling of HIV and Tuberculosis in Kenya, 2012--2017}

\vspace{10pt}
{\normalsize
  Verrah A. Otiende\textsuperscript{1},\;
  Thomas N. Achia\textsuperscript{2},\;
  Henry G. Mwambi\textsuperscript{2}}

\vspace{4pt}
{\small
  \textsuperscript{1}Pan African University PAUSTI, Nairobi, Kenya\quad
  \textsuperscript{2}School of Mathematics, Statistics \& Computer Science,
  University of KwaZulu-Natal, South Africa}

\vspace{4pt}
{\small\itshape
  Corresponding author: verrahodhiambo@gmail.com\quad
  ORCID: 0000-0001-6147-3547}

\vspace{6pt}
{\small
  \textit{Published:} PLOS ONE 15(7): e0234456 (2020)\quad
  DOI: \href{https://doi.org/10.1371/journal.pone.0234456}%
  {\texttt{10.1371/journal.pone.0234456}}}

\vspace{2pt}
{\small
  \textit{Data archive:}
  \href{https://doi.org/10.5281/zenodo.19674273}%
  {\texttt{doi:10.5281/zenodo.19674273}}}
\end{center}

\vspace{8pt}\hrule\vspace{6pt}

%% ── ABSTRACT ─────────────────────────────────────────────────────────────────
\small\noindent\textbf{\textit{Abstract}---}%
This codebook documents all variables, data sources, derived measures, and
analytical indices used in a Bayesian spatio-temporal joint disease modelling
study of HIV and Tuberculosis (TB) co-morbidity across Kenya's 47 counties
from 2012 to 2017. Two CSV files (47 counties $\times$ 16 columns each) contain
annual case notification counts and projected populations for both diseases. The
analytical pipeline computes Standardised Incidence Ratios via indirect
standardisation, fits a shared-component Poisson model in R-INLA, and extracts
spatial, temporal, and space-time interaction effects. The data and code support
full reproduction of all results in the companion publication and can be extended
to other disease pairs or sub-national units in Africa.

\vspace{4pt}
\noindent\textbf{\textit{Keywords}---}%
Bayesian hierarchical model; spatio-temporal analysis; HIV; Tuberculosis;
joint disease modelling; Kenya; R-INLA; shared-component model;
Standardised Incidence Ratio; disease mapping.

\vspace{6pt}\hrule\vspace{6pt}

\small  %% Reduce body/narrative font size throughout

%% ─────────────────────────────────────────────────────────────────────────────
\section{Study Overview}

Table~1 summarises the study design, geographic and temporal scope, the two
diseases modelled, and the statistical framework used throughout the analysis.

\begin{table}[H]
\caption{Study design summary for the HIV and TB joint spatio-temporal analysis.}
\renewcommand{\arraystretch}{1.10}
\scriptsize\begin{tabularx}{\textwidth}{@{}L{3.5cm}X@{}}
\toprule
\textit{Attribute} & \textit{Detail}\\
\midrule
Study period        & 2012--2017 (6 annual time points)\\
Geographic unit     & County ($N = 47$), Kenya\\
Diseases modelled   & HIV, Tuberculosis (TB)\\
Model framework     & Bayesian INLA, shared-component Poisson model\\
Spatial structure   & Besag intrinsic CAR on queen-contiguity adjacency graph\\
Temporal structure  & First-order random walk over 6 years\\
Software            & R ($\geq$ 4.0), R-INLA (\url{www.r-inla.org})\\
\bottomrule
\end{tabularx}
\renewcommand{\arraystretch}{1.10}
\end{table}

%% ─────────────────────────────────────────────────────────────────────────────
\section{File Inventory}

Table~2 lists all supplementary files forming the self-contained research
compendium. Together they provide raw data, a documented analysis script,
a codebook, and spatial boundary data sufficient to reproduce all results
in the companion publication.

\begin{table}[H]
\caption{All supplementary files included in the research compendium.}
\renewcommand{\arraystretch}{1.10}
\scriptsize\begin{tabularx}{\textwidth}{@{}L{5.0cm}L{1.8cm}X@{}}
\toprule
\textit{Filename} & \textit{Format} & \textit{Contents}\\
\midrule
\code{HIV\_Data\_2018.csv}             & CSV        & HIV case counts and projected populations by county, 2012--2017\\
\code{TB\_Data\_2018.csv}              & CSV        & TB notification counts and projected populations by county, 2012--2017\\
\code{population\_by\_county.docx}     & Word       & Annual county population projections from KNBS\\
\code{county\_codes.docx}              & Word       & County codes 1--47 mapped to official county names\\
\code{HIV\_TB\_Kenya\_Analysis.Rmd}    & R Markdown & Documented analysis script with narrative and output explanations\\
\code{HIV\_TB\_Kenya\_Codebook.pdf}    & PDF        & This codebook in journal article format\\
\code{County\_Kenya/County.shp}        & Shapefile  & Kenya county boundary polygons (47 counties, WGS84)\\
\bottomrule
\end{tabularx}
\renewcommand{\arraystretch}{1.10}
\end{table}

%% ─────────────────────────────────────────────────────────────────────────────
\section{TB Dataset Variables (\texttt{TB\_Data\_2018.csv})}

Table~3 describes all 16 columns in the TB dataset. Each row represents one
of Kenya's 47 counties. Columns alternate between annual TB case counts
(\code{TB\_YYYY}) and KNBS projected population denominators (\code{PP\_YYYY})
for each year from 2012 to 2017. The \code{TB\_Total} and \code{PP\_Total}
columns aggregate values across the six-year period and serve as the overall
SIR denominator. \textit{Note on row ordering:} rows are sorted alphabetically
by county name in the raw CSV, not by \code{County\_ID}. The R analysis script
uses \code{match()} to reorder rows to match the shapefile before any spatial
merge.

\begin{table}[H]
\caption{Variable descriptions for \texttt{TB\_Data\_2018.csv} ($N=47$ counties, 16 columns).}
\renewcommand{\arraystretch}{1.10}
\scriptsize\begin{tabularx}{\textwidth}{@{}L{2.7cm}L{1.4cm}L{1.6cm}X@{}}
\toprule
\textit{Variable} & \textit{Type} & \textit{Range} & \textit{Description}\\
\midrule
\code{County\_ID} & Integer & 1--47    & Unique county code. Links to \code{county\_codes.docx} and the INLA spatial index \code{sp.idx}.\\
\code{County}     & String  & ---      & Official county name (KNBS). Rows sorted alphabetically by this field.\\
\code{TB\_2012}   & Integer & $\geq 0$ & New TB cases notified across all NLTP-registered facilities in 2012.\\
\code{PP\_2012}   & Integer & $> 0$    & KNBS projected mid-year population for 2012; the annual person-time denominator.\\
\code{TB\_2013}   & Integer & $\geq 0$ & New TB cases notified in 2013.\\
\code{PP\_2013}   & Integer & $> 0$    & Projected mid-year population for 2013.\\
\code{TB\_2014}   & Integer & $\geq 0$ & New TB cases notified in 2014.\\
\code{PP\_2014}   & Integer & $> 0$    & Projected mid-year population for 2014.\\
\code{TB\_2015}   & Integer & $\geq 0$ & New TB cases notified in 2015.\\
\code{PP\_2015}   & Integer & $> 0$    & Projected mid-year population for 2015.\\
\code{TB\_2016}   & Integer & $\geq 0$ & New TB cases notified in 2016.\\
\code{PP\_2016}   & Integer & $> 0$    & Projected mid-year population for 2016.\\
\code{TB\_2017}   & Integer & $\geq 0$ & New TB cases notified in 2017.\\
\code{PP\_2017}   & Integer & $> 0$    & Projected mid-year population for 2017.\\
\code{TB\_Total}  & Integer & $\geq 0$ & Aggregate TB case count across all six study years (2012--2017).\\
\code{PP\_Total}  & Integer & $> 0$    & Sum of six annual projected populations. Represents total person-years; used as the pooled SIR denominator.\\
\bottomrule
\end{tabularx}
\renewcommand{\arraystretch}{1.10}
\end{table}

%% ─────────────────────────────────────────────────────────────────────────────
\section{HIV Dataset Variables (\texttt{HIV\_Data\_2018.csv})}

Table~4 describes all 16 columns in the HIV dataset. The structure mirrors the
TB file, with case counts (\code{HIV\_YYYY}) and projected populations
(\code{PP\_YYYY}) paired for each year. The key structural difference is the
county identifier column name: the HIV file uses \code{Sn.no} rather than
\code{County\_ID}. Both hold the same integer codes (1--47) and are fully
interchangeable in data joins. All population columns are numerically identical
across both files, as both use the same KNBS inter-censal projection series.

\begin{table}[H]
\caption{Variable descriptions for \texttt{HIV\_Data\_2018.csv} ($N=47$ counties, 16 columns).}
\renewcommand{\arraystretch}{1.10}
\scriptsize\begin{tabularx}{\textwidth}{@{}L{2.7cm}L{1.4cm}L{1.6cm}X@{}}
\toprule
\textit{Variable} & \textit{Type} & \textit{Range} & \textit{Description}\\
\midrule
\code{Sn.no}      & Integer & 1--47    & County serial number. Equivalent to \code{County\_ID} in the TB file; naming difference preserved from original supplementary materials.\\
\code{County}     & String  & ---      & Official county name (KNBS).\\
\code{HIV\_2012}  & Integer & $\geq 0$ & New HIV diagnoses reported to NASCOP surveillance in 2012.\\
\code{PP\_2012}   & Integer & $> 0$    & KNBS projected mid-year population for 2012.\\
\code{HIV\_2013}  & Integer & $\geq 0$ & New HIV diagnoses in 2013.\\
\code{PP\_2013}   & Integer & $> 0$    & Projected mid-year population for 2013.\\
\code{HIV\_2014}  & Integer & $\geq 0$ & New HIV diagnoses in 2014.\\
\code{PP\_2014}   & Integer & $> 0$    & Projected mid-year population for 2014.\\
\code{HIV\_2015}  & Integer & $\geq 0$ & New HIV diagnoses in 2015.\\
\code{PP\_2015}   & Integer & $> 0$    & Projected mid-year population for 2015.\\
\code{HIV\_2016}  & Integer & $\geq 0$ & New HIV diagnoses in 2016.\\
\code{PP\_2016}   & Integer & $> 0$    & Projected mid-year population for 2016.\\
\code{HIV\_2017}  & Integer & $\geq 0$ & New HIV diagnoses in 2017.\\
\code{PP\_2017}   & Integer & $> 0$    & Projected mid-year population for 2017.\\
\code{HIV\_Total} & Integer & $\geq 0$ & Aggregate HIV case count across all six study years.\\
\code{PP\_Total}  & Integer & $> 0$    & Sum of six annual projected populations. Identical in value to \code{PP\_Total} in the TB dataset.\\
\bottomrule
\end{tabularx}
\renewcommand{\arraystretch}{1.10}
\end{table}

%% ─────────────────────────────────────────────────────────────────────────────
\section{County Code Reference}

Table~5 maps all 47 county codes to their official names. These codes appear
throughout the CSV datasets, the INLA model index variables (\code{sp.idx}),
and the posterior result matrices from the fitted model. After the
\code{match()} alignment in the R script, they also correspond to row positions
in the shapefile attribute table (field: \code{COUNTY}).

\begin{table}[H]
\caption{Complete county code reference for all 47 official counties of Kenya.}
\renewcommand{\arraystretch}{1.10}
\scriptsize\begin{tabularx}{\textwidth}{@{}
  C{0.9cm}L{3.9cm}
  C{0.9cm}L{3.9cm}
  C{0.9cm}L{3.7cm}@{}}
\toprule
\textit{ID} & \textit{County} &
\textit{ID} & \textit{County} &
\textit{ID} & \textit{County}\\
\midrule
1  & Turkana        & 17 & Meru        & 33 & Nyamira\\
2  & Marsabit       & 18 & Nandi       & 34 & Narok\\
3  & Mandera        & 19 & Siaya       & 35 & Kisii\\
4  & Wajir          & 20 & Nakuru      & 36 & Murang'a\\
5  & West Pokot     & 21 & Vihiga      & 37 & Migori\\
6  & Samburu        & 22 & Nyandarua   & 38 & Kiambu\\
7  & Isiolo         & 23 & Tharaka     & 39 & Machakos\\
8  & Baringo        & 24 & Kericho     & 40 & Kajiado\\
9  & Keiyo-Marakwet & 25 & Kisumu      & 41 & Nairobi\\
10 & Trans Nzoia    & 26 & Nyeri       & 42 & Makueni\\
11 & Bungoma        & 27 & Tana River  & 43 & Lamu\\
12 & Garissa        & 28 & Kitui       & 44 & Kilifi\\
13 & Uasin Gishu    & 29 & Kirinyaga   & 45 & Taita Taveta\\
14 & Kakamega       & 30 & Embu        & 46 & Kwale\\
15 & Laikipia       & 31 & Homa Bay    & 47 & Mombasa\\
16 & Busia          & 32 & Bomet       &         &\\
\bottomrule
\end{tabularx}
\renewcommand{\arraystretch}{1.10}
\end{table}

%% ─────────────────────────────────────────────────────────────────────────────
\section{Derived Variables and Intermediate Calculations}

All variables in this section are computed at runtime in
\code{HIV\_TB\_Kenya\_Analysis.Rmd} and are not stored in the CSV files.
They are intermediate products of the \textbf{indirect standardisation
procedure}, which adjusts for population size differences before computing
Standardised Incidence Ratios (SIRs). The core formula is
$\mathrm{SIR}_{iy} = O_{iy}/E_{iy}$, where $O_{iy}$ comes from the CSV
and $E_{iy} = \bar{\lambda}_y \times \bar{P}_i$ is derived from the
national crude rate and each county's average annual population.

\subsection{Standard Populations and National Crude Rates}

Table~6 documents the two inputs to the expected count formula. The standard
population $\bar{P}_i = \code{PP\_Total}/6$ is each county's average annual
population over the study period, ensuring the denominator is not biased
towards any single year. Crude rates are computed nationally and serve as the
benchmark: SIR\,$= 1$ means exactly the national rate; SIR\,$> 1$ indicates
excess burden; SIR\,$< 1$ indicates below-average burden.

\begin{table}[H]
\caption{Standard populations and national crude rates used as inputs to the expected count formula $E_{iy} = \bar{\lambda}_y \times \bar{P}_i$.}
\renewcommand{\arraystretch}{1.10}
\scriptsize\begin{tabularx}{\textwidth}{@{}L{3.8cm}L{5.2cm}X@{}}
\toprule
\textit{Variable} & \textit{Formula} & \textit{Description}\\
\midrule
\code{standard.HIV}
  & \code{PP\_Total\,/\,6}
  & Average annual population per county from the HIV file. A 47-element numeric vector.\\
\code{standard.TB}
  & \code{PP\_Total\,/\,6}
  & Same calculation from the TB file. Numerically identical to \code{standard.HIV}.\\[5pt]
\code{cruderate.HIV[y]}
  & $\displaystyle\frac{\sum_i\mathrm{HIV}_{iy}}{\sum_i\mathrm{PP}_{iy}}$
  & National HIV crude rate for year $y\in\{2012,\ldots,2017\}$, in cases per person per year.\\[8pt]
\code{cruderate.HIV\_Total}
  & $\displaystyle\frac{\sum_i\mathrm{HIV\_Total}_i}{\sum_i\mathrm{PP\_Total}_i}$
  & Pooled six-year national HIV crude rate.\\[8pt]
\code{cruderate.TB[y]}
  & $\displaystyle\frac{\sum_i\mathrm{TB}_{iy}}{\sum_i\mathrm{PP}_{iy}}$
  & National TB crude rate for year $y$.\\[8pt]
\code{cruderate.TB\_Total}
  & $\displaystyle\frac{\sum_i\mathrm{TB\_Total}_i}{\sum_i\mathrm{PP\_Total}_i}$
  & Pooled six-year national TB crude rate.\\
\bottomrule
\end{tabularx}
\renewcommand{\arraystretch}{1.10}
\end{table}

\subsection{Expected Counts and Standardised Incidence Ratios}

Table~7 documents the expected case counts and SIRs produced by the
standardisation procedure. Raw SIRs can be statistically unstable for
small-population counties (e.g.\ Lamu, Isiolo) due to small numerators;
the INLA model addresses this by borrowing information from spatial
neighbours through the Besag CAR prior. These continuous SIR values are
later classified into four categories for choropleth mapping (Table~8).

\begin{table}[H]
\caption{Expected case counts and Standardised Incidence Ratios (SIRs) produced by indirect standardisation. Matrices are $47\times6$; vectors have length~47.}
\renewcommand{\arraystretch}{1.10}
\scriptsize\begin{tabularx}{\textwidth}{@{}L{3.8cm}L{5.2cm}X@{}}
\toprule
\textit{Variable} & \textit{Formula} & \textit{Description}\\
\midrule
\code{expected.HIV[,y]}
  & $\bar{\lambda}^{\mathrm{HIV}}_y\times\bar{P}^{\mathrm{HIV}}_i$
  & Expected HIV cases for county $i$, year $y$; $47\times6$ matrix.\\
\code{expected.HIV\_Total}
  & $\bar{\lambda}^{\mathrm{HIV}}_\bullet\times\bar{P}^{\mathrm{HIV}}_i$
  & Pooled expected HIV count; 47-element vector.\\
\code{expected.TB[,y]}
  & $\bar{\lambda}^{\mathrm{TB}}_y\times\bar{P}^{\mathrm{TB}}_i$
  & Expected TB cases; $47\times6$ matrix.\\
\code{expected.TB\_Total}
  & $\bar{\lambda}^{\mathrm{TB}}_\bullet\times\bar{P}^{\mathrm{TB}}_i$
  & Pooled expected TB; 47-element vector.\\[4pt]
\code{adj.HIV[,y]}
  & $O^{\mathrm{HIV}}_{iy}/E^{\mathrm{HIV}}_{iy}$
  & Continuous HIV SIR per county-year. Values $>1$ indicate above-national-average burden; $47\times6$ matrix.\\
\code{adj.HIV.total}
  & $O^{\mathrm{HIV}}_{\bullet i}/E^{\mathrm{HIV}}_{\bullet i}$
  & Pooled HIV SIR across all six years; 47-element vector.\\
\code{adj.TB[,y]}
  & $O^{\mathrm{TB}}_{iy}/E^{\mathrm{TB}}_{iy}$
  & Continuous TB SIR; $47\times6$ matrix.\\
\code{adj.TB.total}
  & $O^{\mathrm{TB}}_{\bullet i}/E^{\mathrm{TB}}_{\bullet i}$
  & Pooled TB SIR; 47-element vector.\\
\bottomrule
\end{tabularx}
\renewcommand{\arraystretch}{1.10}
\end{table}

\subsection{SIR Categories for Choropleth Mapping}

Table~8 defines the four ordered categories used to discretise continuous
SIRs for colour-coded choropleth mapping. The same breakpoints are applied
uniformly across all six years and both diseases so the colour scale is
directly comparable across all map panels. Maps use the \textbf{YlOrRd}
(Yellow-Orange-Red) ColorBrewer palette; space-time exceedance maps use
\textbf{YlGn} (Yellow-Green) to visually distinguish evidence maps from
risk-magnitude maps.

\begin{table}[H]
\caption{SIR classification scheme applied uniformly across all years and both diseases. Breakpoints: $c=(0,\;0.40,\;0.70,\;0.90,\;6)$.}
\renewcommand{\arraystretch}{1.10}
\scriptsize\begin{tabularx}{\textwidth}{@{}L{2.8cm}L{2.8cm}X@{}}
\toprule
\textit{Category} & \textit{SIR range} & \textit{Epidemiological interpretation}\\
\midrule
\code{0.00--0.40} & $[0,\;0.40]$    & Very low burden. County incidence rate is below 40\% of the national average. Mapped in pale yellow.\\
\code{0.41--0.70} & $(0.40,\;0.70]$ & Low burden. County rate between 40\% and 70\% of the national average.\\
\code{0.71--0.99} & $(0.70,\;0.90]$ & Moderate burden. Approaching but still below the national average.\\
\code{1.00+}      & $>0.90$         & High burden. At or above the national average. Mapped in red; indicates priority area for public health intervention.\\
\bottomrule
\end{tabularx}
\renewcommand{\arraystretch}{1.10}
\end{table}

%% ─────────────────────────────────────────────────────────────────────────────
\section{INLA Model: Data Structure and Index Variables}

The joint Bayesian model uses a long-format data frame with
$47\times6\times2 = 564$ rows, one per county-year-disease combination.
The first 282 rows correspond to HIV (47 counties $\times$ 6 years);
rows 283--564 correspond to TB in the same county-year order.

\subsection{Long-Format Data Frame Variables}

Table~9 lists all columns in the 564-row data frame passed to \code{inla()}.
The \code{Expected} column serves as the Poisson offset so the model estimates
a log relative risk rather than an absolute rate. The \code{intercept} factor
provides disease-specific intercepts via the \code{$-1+$intercept} formula syntax.

\begin{table}[H]
\caption{Variables in the 564-row long-format INLA data frame. Rows 1--282 are HIV observations; rows 283--564 are TB observations.}
\renewcommand{\arraystretch}{1.10}
\scriptsize\begin{tabularx}{\textwidth}{@{}L{3.0cm}L{1.6cm}L{2.0cm}X@{}}
\toprule
\textit{Variable} & \textit{Type} & \textit{Range} & \textit{Description}\\
\midrule
\code{County.ID}      & Integer & 1--47      & County identifier, recycled across years and diseases.\\
\code{Year}           & Integer & 2012--2017 & Calendar year of observation.\\
\code{County.Year.ID} & Integer & 1--282     & Sequential county$\times$year index.\\
\code{Expected}       & Numeric & $>0$       & Expected count; the Poisson offset $E$ in \code{inla()}.\\
\code{Observed}       & Integer & $\geq 0$   & Observed new case count per county-year. The response variable.\\
\code{Disease}        & String  & HIV / TB   & Labels each row as HIV or TB.\\
\code{sp.idx}         & Integer & 1--47      & Spatial county index for Besag CAR random effects.\\
\code{tm.idx}         & Integer & 1--6       & Temporal index: Year $-$ 2011 ($2012\mapsto1,\,\ldots,\,2017\mapsto6$).\\
\code{sp.tm.idx}      & Integer & 1--282     & Combined space-time interaction index.\\
\code{dis.idx}        & Integer & 1 or 2     & Disease indicator: 1 = HIV, 2 = TB.\\
\code{intercept}      & Factor  & 1 or 2     & Factor version of \code{dis.idx}; provides disease-specific intercepts.\\
\bottomrule
\end{tabularx}
\renewcommand{\arraystretch}{1.10}
\end{table}

\subsection{Disease-Specific Index Variables (NA Structure)}

Table~10 describes the NA-structured index columns that define the joint model.
For a disease-specific effect, the index is active only for that disease's rows
and \texttt{NA} elsewhere. INLA treats NA as a non-contributing observation.
Shared fields use a single index active for all 564 rows; each disease enters
the shared field via a \textit{copy} term with an estimated positive scaling
coefficient $\beta$.

\begin{table}[H]
\caption{Disease-specific and shared index variables using the NA structure. An \texttt{NA} entry tells INLA the row does not contribute to that random effect.}
\renewcommand{\arraystretch}{1.10}
\scriptsize\begin{tabularx}{\textwidth}{@{}L{3.3cm}L{2.8cm}X@{}}
\toprule
\textit{Variable} & \textit{Active for} & \textit{Description}\\
\midrule
\code{sp.idx1\,/\,sp.1} & HIV rows only  & HIV spatial index: county codes 1--47 for HIV; \texttt{NA} for all 282 TB rows.\\
\code{sp.idx2\,/\,sp.2} & TB rows only   & TB spatial index; \texttt{NA} for all HIV rows.\\
\code{sp.dum}            & Both diseases  & Shared spatial field; both diseases scaled by estimated $\beta^{sp}$.\\
\code{tm.idx1\,/\,tm.1} & HIV rows only  & HIV temporal index (1--6); \texttt{NA} for TB rows.\\
\code{tm.idx2\,/\,tm.2} & TB rows only   & TB temporal index; \texttt{NA} for HIV rows.\\
\code{tm.dum}            & Both diseases  & Shared temporal field.\\
\code{sp.tm.idx1}        & HIV rows only  & HIV space-time interaction index; \texttt{NA} for TB rows.\\
\code{sp.tm.idx2}        & TB rows only   & TB space-time interaction index; \texttt{NA} for HIV rows.\\
\code{sp.tm.dum}         & Both diseases  & Shared space-time interaction field; active for all 564 rows.\\
\bottomrule
\end{tabularx}
\renewcommand{\arraystretch}{1.10}
\end{table}

\subsection{Adjacency Structures}

Table~11 describes the two adjacency matrices passed to the \code{graph}
argument of each \code{f()} term. Together they determine how information
is borrowed across space and time during Bayesian smoothing, stabilising
estimates for small-population county-year cells.

\begin{table}[H]
\caption{Spatial and temporal adjacency matrices used as graph structures in the INLA model formula.}
\renewcommand{\arraystretch}{1.10}
\scriptsize\begin{tabularx}{\textwidth}{@{}L{2.8cm}L{1.8cm}X@{}}
\toprule
\textit{Object} & \textit{Dimensions} & \textit{Description}\\
\midrule
\code{w.sp} & $47\times47$ & Binary queen-contiguity matrix. Entry $(i,j)=1$ if counties $i$ and $j$ share any boundary point or vertex; derived via \code{poly2nb()} and \code{nb2mat()}. Average neighbours per county: $\approx5.5$.\\
\code{w.tm} & $6\times6$   & First-order random walk matrix. Each year connected only to its immediate predecessor and successor; encodes smooth temporal dependency.\\
\code{kenya.graph} & Text file & INLA adjacency file written by \code{nb2INLA()}; passed to the \code{graph} argument of each Besag CAR \code{f()} term.\\
\bottomrule
\end{tabularx}
\renewcommand{\arraystretch}{1.10}
\end{table}

\subsection{Prior Specifications}

Table~12 summarises the two prior distributions. Both are weakly informative,
permitting a broad range of plausible parameter values. The Normal prior on
copy-effect $\beta$ is constrained to $(0,\infty)$, reflecting the known
positive HIV--TB epidemiological linkage.

\begin{table}[H]
\caption{Prior distributions specified in the INLA model.}
\renewcommand{\arraystretch}{1.10}
\scriptsize\begin{tabularx}{\textwidth}{@{}L{3.2cm}L{4.5cm}X@{}}
\toprule
\textit{Prior} & \textit{Parameters} & \textit{Applied to}\\
\midrule
Log-Gamma on precision
  & shape $=0.5$, rate $=0.0005$
  & All spatial (Besag), temporal (Besag), and space-time (IID) precision hyperparameters.\\
Normal on copy $\beta$
  & $\mu=0$, $\tau=1/5.9$, SD$\approx2.43$, range $(0,\infty)$
  & Scaling coefficients for shared spatial (\code{prior.beta.s}), temporal (\code{prior.beta.t}), and space-time (\code{prior.beta.s.t}) copy effects.\\
\bottomrule
\end{tabularx}
\renewcommand{\arraystretch}{1.10}
\end{table}

%% ─────────────────────────────────────────────────────────────────────────────
\section{INLA Model Output Variables}

The fitted model object stores posterior summaries in
\code{model.inla\$summary.random}. Results are extracted by exponentiating
posterior means to convert from log-RR to the natural scale, and by computing
posterior exceedance probabilities.

\subsection{Spatial Random Effects}

Table~13 describes five posterior spatial relative risk (RR) variables
appended to the spatial data frame \code{kenya.one} for choropleth mapping.
Three types are reported: the \textbf{shared component} (the co-epidemic
spatial signal); the \textbf{disease-specific residual} after removing the
shared pattern; and the \textbf{combined total} — the most policy-relevant
output, integrating both into a full spatial burden estimate per disease.
Categorical versions use breakpoints $(0,\;0.40,\;0.70,\;1.00,\;6)$.

\begin{table}[H]
\caption{Posterior mean spatial relative risk variables from the INLA model (exponentiated, natural scale). Categorical versions use breakpoints $(0,\;0.40,\;0.70,\;1.00,\;6)$.}
\renewcommand{\arraystretch}{1.10}
\scriptsize\begin{tabularx}{\textwidth}{@{}L{3.2cm}L{5.0cm}X@{}}
\toprule
\textit{Variable} & \textit{Formula} & \textit{Description}\\
\midrule
\code{SHARED\_C}
  & $\exp(\hat{\phi}_i)$
  & Shared spatial RR. Values $>1$ indicate both diseases are jointly elevated in county $i$. The co-epidemic spatial signal.\\
\code{HIV\_C}
  & $\exp(\hat{u}_{\mathrm{HIV},i})$
  & HIV-specific residual spatial RR after removing the shared component.\\
\code{TB\_C}
  & $\exp(\hat{u}_{\mathrm{TB},i})$
  & TB-specific residual spatial RR.\\
\code{Total.HIV\_C}
  & $\exp\!\bigl(\hat{u}_{\mathrm{HIV},i}+\hat{\beta}^{sp}_{\mathrm{HIV}}\hat{\phi}_i\bigr)$
  & Combined HIV spatial RR (disease-specific + scaled shared). Primary HIV spatial output in the companion article.\\
\code{Total.TB\_C}
  & $\exp\!\bigl(\hat{u}_{\mathrm{TB},i}+\hat{\beta}^{sp}_{\mathrm{TB}}\hat{\phi}_i\bigr)$
  & Combined TB spatial RR. Primary TB spatial output.\\
\bottomrule
\end{tabularx}
\renewcommand{\arraystretch}{1.10}
\end{table}

\subsection{Space-Time Interaction and Exceedance Probabilities}

Table~14 describes the shared space-time interaction variables. These capture
residual county-year heterogeneity beyond the main spatial and temporal
patterns. The exceedance probability $P(\delta_{it}>0\mid\mathbf{y})$ integrates
the full posterior distribution and is more reliable than the posterior mean
alone. Values $>0.75$ provide very strong evidence of excess shared risk in that
county-year combination.

\begin{table}[H]
\caption{Shared space-time interaction variables from the INLA model. Computed for all $47\times6=282$ county-year combinations.}
\renewcommand{\arraystretch}{1.10}
\scriptsize\begin{tabularx}{\textwidth}{@{}L{3.2cm}X@{}}
\toprule
\textit{Variable} & \textit{Description}\\
\midrule
\code{delta.matrix\_s}
  & $47\times6$ matrix of posterior mean shared space-time effects on the log scale. Rows = counties (shapefile order); columns = years 2012--2017. Values $>0$ indicate excess shared risk beyond the main spatial and temporal components.\\
\code{pp.delta\_s}
  & Exceedance probability $P(\delta_{it}>0\mid\mathbf{y})$ per county-year, via \code{inla.pmarginal()}. Values near 1 = very strong evidence of excess risk; values near 0 = evidence of below-expected risk.\\
\code{pp.delta.factor\_s}
  & Categorical exceedance probability in four bands: \textbf{0.00--0.25} (below-expected); \textbf{0.26--0.50} (little evidence of excess); \textbf{0.51--0.75} (moderate evidence); \textbf{0.76--1.00} (very strong evidence of excess shared risk). Displayed with \textbf{YlGn} palette.\\
\bottomrule
\end{tabularx}
\renewcommand{\arraystretch}{1.10}
\end{table}

%% ─────────────────────────────────────────────────────────────────────────────
\section{Data Notes and Caveats}

Table~15 flags known data limitations, naming inconsistencies, and
implementation choices that analysts should be aware of before using or
extending this analysis.

\begin{table}[H]
\caption{Known data limitations, naming inconsistencies, and implementation notes.}
\renewcommand{\arraystretch}{1.10}
\scriptsize\begin{tabularx}{\textwidth}{@{}L{4.0cm}X@{}}
\toprule
\textit{Issue} & \textit{Detail}\\
\midrule
Row ordering
  & CSV files are sorted alphabetically by county name, not by \code{County\_ID}. Use \code{match()} on the shapefile \code{COUNTY} field to reorder rows before any spatial merge. Skipping this step silently assigns SIR values to incorrect county polygons on the map.\\
Column name discrepancy
  & The HIV file uses \code{Sn.no}; the TB file uses \code{County\_ID}. Both hold codes 1--47 and are interchangeable. The inconsistency is preserved from the original PLOS ONE supplementary files.\\
Pooled denominator
  & \code{PP\_Total} represents total person-years over six years, not a single-year count. Do not use it as a point-in-time population denominator.\\
INLA temporal indexing
  & All indices must start at 1. Use $\code{tm.idx}=\text{Year}-2011$: $2012\mapsto1$, \ldots, $2017\mapsto6$.\\
Copy-effect $\beta$
  & Range constrained to $(0,\infty)$. $\beta\approx0$: disease-specific effects dominate. $\beta\approx1$: disease closely follows the shared pattern.\\
Reproducibility
  & INLA model fitting: $\approx$20--40 minutes. Use \code{cache=TRUE} to avoid re-running on subsequent knits. Fitted model saved to \code{st\_model\_GAMMA.RData}.\\
\bottomrule
\end{tabularx}
\renewcommand{\arraystretch}{1.10}
\end{table}

\vspace{12pt}\hrule\vspace{6pt}

\noindent\small\textit{End of Codebook.}\enspace
Otiende VA, Achia TN, Mwambi HG (2020).
Bayesian hierarchical modeling of joint spatiotemporal risk patterns for HIV
and TB in Kenya. \textit{PLOS ONE} 15(7): e0234456.
DOI: \texttt{10.1371/journal.pone.0234456}

\end{document}
